    \begin{spacing}{1.5}
        
        \begin{center}
            {\large\bfseries INTRODUCTION}
        \end{center}
        \addcontentsline{toc}{chapter}{INTRODUCTION}
        Depuis la création de l'internet, le commerce en ligne a connu une forte croissance et participe désormais 
        de manière significative à l'économie mondiale. 
        Les plateformes e-commerce traditionnelles reposent souvent sur des architectures monolithiques, 
        où toutes les fonctionnalités sont regroupées dans une seule application.\\
        Bien que cette approche ait permis un développement rapide, elle présente plusieurs limitations, 
        notamment en termes de scalabilité, de maintenabilité et d'agilité de développement.
        C'est pourquoi de nombreuses entreprises se tournent vers des architectures basées sur des microservices, 
        qui offrent une plus grande flexibilité et une meilleure capacité à s'adapter aux évolutions du marché.
        Cette transition vers les microservices nécessite une réflexion approfondie sur l'architecture logicielle, 
        les choix technologiques et les processus de développement.\\
        De cette analyse, nous avons dégagé une problématique centrale : comment concevoir une plateforme e-commerce
        basée sur des microservices tout en garantissant une expérience utilisateur fluide et cohérente ? \\
        Pour répondre à cette problématique, l'objectif de ce mémoire est de migrer une architecture monolithique
        existante vers une architecture basée sur des microservices en utilisant des methodologies et de techniques modernes 
        de développement logiciel.
        Nous allons explorer les différentes étapes de cette migration,
        en mettant l'accent sur les aspects techniques, organisationnels et stratégiques.
        Nous aborderons également les défis rencontrés lors de cette transition et les solutions mises en place pour les surmonter.\\
        Ce mémoire est structuré en plusieurs chapitres.
        Le premier chapitre présente le contexte tel que la présentation des parties prenantes et le cadre du projet.
        Le second chapitre analyse l'architecture existante et la conception de l'architecture cible.
        Le troisième chapitre détaille le processus de migration, les outils utilisés et les bonnes pratiques adoptées.
        Enfin, le dernier chapitre évalue les résultats obtenus et propose des pistes d'amélioration pour l'avenir.
    \end{spacing}
